% Copyright Myeongjin Kang. Original prose: CC BY 4.0.
\documentclass[11pt]{article}
\usepackage[a4paper,margin=29mm]{geometry}
\usepackage{amsmath,amssymb,amsthm,mathtools}
\usepackage{microtype}
\usepackage[T1]{fontenc}
\usepackage{lmodern}
\usepackage{enumitem}
\usepackage[colorlinks=true,linkcolor=blue!50!black,citecolor=blue!50!black,urlcolor=blue!50!black]{hyperref}
\usepackage{xcolor}
\hypersetup{pdftitle={One integral mapping-class orbit in every rank},pdfauthor={Myeongjin Kang}}
\newtheorem{theorem}{Theorem}[section]
\newtheorem{proposition}[theorem]{Proposition}
\newtheorem{lemma}[theorem]{Lemma}
\newtheorem{corollary}[theorem]{Corollary}
\theoremstyle{definition}\newtheorem{definition}[theorem]{Definition}
\theoremstyle{remark}\newtheorem{remark}[theorem]{Remark}
\newcommand{\C}{\mathbb C}\newcommand{\Q}{\mathbb Q}\newcommand{\Z}{\mathbb Z}
\newcommand{\OO}{\mathcal O}\newcommand{\F}{\mathcal F}
\newcommand{\SL}{\mathrm{SL}}\newcommand{\GL}{\mathrm{GL}}
\newcommand{\Mod}{\operatorname{Mod}}\newcommand{\Ad}{\operatorname{Ad}}
\newcommand{\tr}{\operatorname{tr}}\newcommand{\diag}{\operatorname{diag}}
\newcommand{\Hom}{\operatorname{Hom}}\newcommand{\im}{\operatorname{im}}
\newcommand{\slr}{\mathfrak{sl}_r}\newcommand{\h}{\mathfrak h}
\newcommand{\hc}{\mathfrak h_C}\newcommand{\TC}{T_C}
\newcommand{\X}{\mathcal X_{g,r}}\newcommand{\R}{\mathcal R_{g,r}}
\newcommand{\dd}{\mathrm d}\newcommand{\eps}{\varepsilon}
\setlength{\emergencystretch}{2em}
\numberwithin{equation}{section}
\title{One integral mapping-class orbit in every rank}
\author{Myeongjin Kang}
\date{October 18, 2026}
\begin{document}
\maketitle
\begin{abstract}
For every closed oriented surface of genus $g\geq2$ and every $r\geq2$,
we construct a representation into $\SL_r(\OO_L)$, for a number field
$L\subset\C$, whose character has Zariski-dense orbit in the entire complex
$\SL_r$ character variety under the orientation-preserving mapping class group.
The rank-two assertion is due to Coccia--Litt and has further antecedents.
For $r\geq3$, two families of simple curves give a full cotangent certificate
at a representation of the form $(A,B,C,I,\ldots,I)$ with $A,B$ diagonal.
Their joint residual kernel is eliminated by a symmetric/skew-symmetric
Cartan degeneration. Arithmetic parameters are chosen by finitely many open
conditions; the countably many nonresonance conditions are imposed only on
a complex torus already contained in one arithmetic orbit closure.
Classical twist Hamiltonians then force that closure to be the whole variety.
The higher-rank construction is presented as new to this project, without
a claim of worldwide priority.
\end{abstract}
\noindent\small Author: no institutional affiliation. Original exposition is licensed
under CC BY 4.0; accompanying original code is licensed under MIT.\normalsize
\tableofcontents
\section{Statement, inputs, and antecedents}\label{sec:context}
Let $\Sigma_g$ be a closed connected oriented surface of genus $g$. Put
\[
 \R=\Hom(\pi_1(\Sigma_g),\SL_r(\C)),\qquad
 \X=\R\mathbin{/\!\!/}\SL_r(\C),
\]
where the latter is the affine character variety. The group $\Mod(\Sigma_g)$
consists of orientation-preserving mapping classes and acts on $\X$ by
precomposition; inner automorphisms of the surface group have no effect on
characters. We use $[x,y]=xyx^{-1}y^{-1}$ throughout.

\begin{theorem}\label{thm:main}
For every pair of integers $g\geq2$ and $r\geq2$, there exist a number field
$L\subset\C$ and a homomorphism
\[
 \rho:\pi_1(\Sigma_g)\longrightarrow\SL_r(\OO_L)
\]
such that $\Mod(\Sigma_g)\cdot[\rho]$ is Zariski dense in $\X$.
\end{theorem}

Here $\OO_L$ is the full ring of integers. The conclusion concerns one
orbit of an actual integral representation and the entire complex character
variety. There is no generic-input condition and no restriction to a marked
locus. No assertion for punctured surfaces is made. For $r\geq3$ the field
in our construction depends only on $r$, although the selected matrices may
depend on $g$.

\subsection{The credited rank-two case}
Coccia--Litt \cite[Propositions 3.1.10 and 5.0.1, Theorem 5.0.2]{CL}
already prove the assertion for $r=2$, in a more general relative setting.
Their Proposition 5.0.1 produces an actual integral representation that is
good for a pants decomposition, and Proposition 3.1.10 proves density of its
mapping-class orbit. To meet the unit hypothesis one can use
$K=\Q(\sqrt5)$ and $v=(1+\sqrt5)/2$: both $v$ and $v-v^{-1}=1$ are units,
and $v$ is not a root of unity. Their field extension and construction
therefore apply to every closed genus $g\geq2$. Density over $\overline\Q$
remains density after extending scalars to $\C$: for each finite-dimensional
space of polynomials of bounded degree, the evaluation conditions have
coefficients in $\overline\Q$, so their kernel commutes with extension of
scalars. The remainder of the proof assumes $r\geq3$.

Litt's survey \cite[Theorem 5.4.13]{Litt} gives another rank-two
single-orbit route, using a Shimura-curve local system with a unitary Galois
conjugate and Previte--Xia's Theorem 1.4 \cite{PreviteXia}.
Golsefidy--Tamam \cite[Section 3.3 and Theorems G--J]{GT}
identify exceptional loci that must be excluded in the general formulation
of that Previte--Xia theorem. We record both the antecedent and this caveat;
we do not use the uncaveated statement as an input to our rank-two proof.
Golsefidy--Tamam's Theorem I gives a substantially stronger rank-two
orbit-closure theorem, under its sufficiently-dense, nonexceptional,
infinite-orbit hypotheses. Their detailed orbit-closure preprint
\cite{GTpre} explicitly treats $\SL_2$.

\subsection{What the higher-rank argument adds}
The pinned manuscript \cite{Source027} proves potential density of the
\emph{whole set} of integral characters in every rank, already over one
number field. Its closed-surface argument first puts a whole fixed-spectrum
fiber in that closure, then uses twist-torus saturation, Baire avoidance,
and a submersion. Those operations are antecedents here; density of a union
of arithmetic orbits does not itself choose one orbit.

Goldman \cite{Goldman} supplies the symplectic and Hamiltonian twist
machinery. Pickrell--Xia \cite[Proposition 2.1.24 and Lemma 2.1.34]{PX}
provide a close antecedent to infinitesimal spanning: centralizers of
products $hg^n$, dense torus powers, and a BCH/root-component argument,
followed by transverse sewing in their Section 3. Their result concerns
compact groups and generic fibers in an almost-everywhere argument.
The task here is to certify a particular complex family with two commuting
handle matrices inside the closure of one arithmetic orbit.

For comparison, the commuting twist tori associated to a pants decomposition
have dimension at most $(3g-3)(r-1)$. For $r\geq3$,
\[
 (3g-3)(r-1)<(g-1)(r^2-1)=\tfrac12\dim\X.
\]
Thus the rank-two pants-torus fiber argument does not simply give the required
higher-rank cotangent certificate. We use two noncommuting families of simple
curves and reduce their common kernel to
\begin{equation}\label{eq:residual-intro}
 (\Ad_A-\Ad_C)S=0,\qquad (I-\Ad_C)S\in\h+\hc.
\end{equation}
A symmetric/skew-symmetric degeneration proves this system generically
injective. Arithmetic parameters are then chosen in finitely many nonempty
opens. A further complex parameter, already inside the arithmetic orbit
closure, is chosen outside countably many resonance sets.

Golsefidy--Tamam \cite[Section 2.4, second remark]{GT} explicitly suggest
extending their results to $\SL_d$ and other Chevalley groups using
Larsen--Pink. That extension is not proved there. Our existence theorem
is a special case of the anticipated broader higher-rank picture, rather
than a claim that such a picture was unanticipated. The construction and
proof are new to this project; bounded searches found no higher-rank
single-orbit integral theorem, but do not establish worldwide novelty.

The choice of point cannot be replaced by an unexamined high-rank
dense-image assertion. A MathOverflow comment by M.~Kapovich reports
non-dense mapping-class orbits of Zariski-dense high-rank unitary
representations \cite{MO}. We cite that as a discussion pointer, not as a
proved theorem used here. Independently, the actual paper of
Biswas--Koberda--Mj--Santharoubane
\cite[Proposition 1.2 and Section 5]{BKMS} constructs infinite-image
representations whose \emph{$\GL_n$ characters} are fixed by the entire
mapping class group, using TQFT representations and point pushing.
That is the precise published statement we use for comparison; it is
not the same as the stronger dense-image observation in the comment.
Neither observation is needed as a lemma in our proof.

\subsection{Geometric input and conventions}\label{subsec:geometry}
We use the following exact input from \cite{Source027}: the determinant-one
fiber
\[
 \{(A_1,B_1,\ldots,A_g,B_g)\in\SL_r^{2g}:
       \prod_{i=1}^g[A_i,B_i]=I\}
\]
is irreducible of dimension $(2g-1)(r^2-1)$. This is the
$\SL_r$ intermediate assertion established \emph{inside the proof} of
Proposition \texttt{prop:surface-irreducibility}, in
\texttt{surfaces.tex}, before the scalar-extension paragraph. The
proposition's stated conclusion concerns the $\GL_r$ fiber; we do not
infer determinant-one irreducibility from that conclusion. The continuous
surjective quotient morphism then makes $\X$ irreducible.

At a smooth irreducible representation the tangent space of $\X$ is the
kernel of the linearized surface relation modulo infinitesimal simultaneous
conjugation. The trace pairing on $\slr$ and cup product give the
nondegenerate holomorphic Goldman form on this smooth locus
\cite[Introduction and Sections 2--3]{Goldman}. We will check smoothness
at our chosen representations directly. All flow arguments take place in
an ambient neighborhood in this smooth irreducible locus, hence in
Goldman's $D/G$; the orbit closure itself need not be smooth.

Write the first two pairs of surface generators as $(a,b),(c,d_0)$ and the
others as $(a_i,b_i)$ for $3\leq i\leq g$. Their matrices are
$(A,B,C,D),(A_i,B_i)$. Let $T\subset\SL_r$ be the diagonal torus and
$\h\subset\slr$ its Lie algebra. Given a regular semisimple $C=QC_\mu Q^{-1}$,
write $\TC=QTQ^{-1}$ and $\hc=Q\h Q^{-1}$.

\section{An arithmetic torus with dense cyclic generators}\label{sec:arithmetic}
The following elementary unit construction fixes the number field before
any genericity condition is imposed.

\begin{lemma}\label{lem:units}
For every $r\geq2$ there are a number field $L\subset\C$ and units
$\mu_1,\ldots,\mu_r\in\OO_L^\times$ with product one such that
\[
 \prod_i\mu_i^{e_i}=1\quad\Longleftrightarrow\quad
 e_1=\cdots=e_r\qquad(e_i\in\Z).
\]
In particular, $C_\mu=\diag(\mu_1,\ldots,\mu_r)$ generates a Zariski-dense
cyclic subgroup of $T$ and has distinct eigenvalues.
\end{lemma}
\begin{proof}
Fix an integer $c_0\geq3$ and set $f(X)=X^r-c_0X^{r-1}-1$.
On $|X|=1$, the middle term has modulus $c_0>2\geq|X^r-1|$.
Rouch\'e's theorem gives exactly $r-1$ roots inside the unit circle.
Since $f(c_0)=-1$ and $f(x)$ tends to infinity on the positive real axis,
there is one real root greater than $c_0$, the unique root outside.
If $f$ were reducible over $\Q$, Gauss's lemma would give a nonconstant
monic integral factor not containing that outside root. Its constant term
would be a nonzero integer of modulus less than one, a contradiction.
Thus $f$ is irreducible; characteristic zero gives distinct roots
$\alpha_1,\ldots,\alpha_r$. They are algebraic units because $f$ is monic
with constant term $-1$.

Suppose $\prod_i\alpha_i^{e_i}$ is a root of unity. Choose an index whose
exponent is $M=\max_i e_i$. Irreducibility allows a Galois automorphism
$\sigma$ of the splitting field that sends this root to the unique outside
root. Since $|\prod_i\sigma\alpha_i|=1$, we have
\[
 0=\sum_i e_i\log|\sigma\alpha_i|
   =\sum_i(e_i-M)\log|\sigma\alpha_i|.
\]
The term for the outside root is zero; all other logarithms are strictly
negative and all their coefficients are nonpositive. Every term is
therefore nonnegative, and equality forces $e_i=M$ for every $i$.

Choose a root of unity $\eta$ with $\eta^r=(-1)^{r-1}$, and put
$\mu_i=\eta\alpha_i$. Since $\prod_i\alpha_i=(-1)^{r+1}$, their product is
one. A relation among the $\mu_i$ gives a root-of-unity relation among the
$\alpha_i$, so its exponents are all equal. Conversely all equal exponents
give a relation. Take $L$ to be the splitting field with $\eta$ adjoined.
The character lattice of $T$ is $\Z^r/\Z(1,\ldots,1)$. A proper algebraic
subgroup of a torus is annihilated by a nontrivial character, so the stated
relation lattice proves cyclic density. Distinctness follows already from
that of the $\alpha_i$.
\end{proof}

For a tuple $m=(m_1,\ldots,m_{r-1})$ of positive integers define
\begin{equation}\label{eq:Bfamily}
 B(m)=\diag\left(\mu_1^{m_1},\ldots,\mu_{r-1}^{m_{r-1}},
                    \prod_{i<r}\mu_i^{-m_i}\right).
\end{equation}
\begin{lemma}\label{lem:arithmetic-dense}
The set of pairs $(B(m),Q)$, with all $m_i\geq1$ and
$Q\in\SL_r(\Z)$, is Zariski dense in $T\times\SL_r(\C)$.
Every $B(m)$ generates $T$ densely. Every
$C=QC_\mu Q^{-1}$ belongs to $\SL_r(\OO_L)$ and generates $\TC$ densely.
\end{lemma}
\begin{proof}
Each $\mu_i$ has infinite order, by Lemma~\ref{lem:units}. The first $r-1$
coordinates in \eqref{eq:Bfamily} range independently over infinite subsets
of $\C^\times$. A Laurent polynomial vanishing on their product is zero,
by induction on the number of variables. This proves density of the
$B(m)$. If a character with exponents $e_i$ is one on $B(m)$, then
\[
 \prod_{i<r}\mu_i^{m_i(e_i-e_r)}=1.
\]
Apply Lemma~\ref{lem:units} with final exponent zero. All exponents are
zero, hence all $e_i=e_r$, proving cyclic density.

The Zariski closure of $\SL_r(\Z)$ contains each elementary root group:
it contains its integral parameters, an infinite subset of an affine line.
Those root groups generate $\SL_r(\C)$, proving density. Products of dense
sets are dense. Finally $Q^{-1}$ has integral entries because $\det Q=1$,
so $QC_\mu Q^{-1}$ is integral; its torus density follows by conjugation.
\end{proof}

All later parameter opens will be taken inside this single fixed-eigenvalue
family. We never assume that a countable set of arithmetic points avoids a
countable union of proper closed sets.

\section{Two families of simple curves}\label{sec:curves}
We fix the based-word convention, since the order of the factors matters
both topologically and in the tangent calculation. A genus-two surface
with one boundary component is obtained by thickening a rose with cyclic
half-edge order
\begin{equation}\label{eq:rose}
 a,\ b^{-1},\ a^{-1},\ b,\ c,\ d_0^{-1},\ c^{-1},\ d_0.
\end{equation}
Composing the cyclic permutation with edge reversal gives one boundary
cycle with word $[a,b][c,d_0]$. In genus $g$ append the blocks
$a_i,b_i^{-1},a_i^{-1},b_i$, and cap the boundary to obtain $\Sigma_g$.
A turn from a letter $x$ to $y$ is drawn as a chord joining half-edges
$x^{-1}$ and $y$ in the vertex disk.

\begin{lemma}\label{lem:curves}
For all $k,l\geq1$, the words
\begin{align}
 u_0(k,l)&=b a^k d_0 c^l,\label{eq:u0}\\
 u_1(k,l)&=a b a^{k-1}c d_0 c^{l-1}\label{eq:u1}
\end{align}
represent simple nonseparating curves. For every $i\geq3$, so do
$u_1(k,l)a_i$ and $u_1(k,l)b_i^{-1}$.
\end{lemma}
\begin{proof}
First consider the words $bd_0$ and $abcd_0$. Their turn chords in
\eqref{eq:rose}, numbered from one to eight, are respectively
\[
 \{2,8\},\{4,6\};\qquad
 \{3,4\},\{2,5\},\{7,8\},\{1,6\}.
\]
In each case the chords are noninterleaving. Each edge is used once,
so these chords and the edge bands smooth to a single embedded closed
curve. Its homology class is nonzero and primitive, hence it is
nonseparating after capping the boundary.

The oriented Dehn twist about $a$, or its inverse according to the choice
of a positive twist, has based action $b\mapsto ba$ and fixes $a$ and the
other generators. This follows by cutting the $a$-annulus and following the
single transverse $b$-arc across one twist. Iteration gives
$b\mapsto ba^m$. The identical construction on the other handle gives
$d_0\mapsto d_0c^n$. These are actual boundary-fixing twist powers;
algebraically $[a,ba^m]=[a,b]$ and $[c,d_0c^n]=[c,d_0]$, so the chosen
based boundary word is fixed exactly. Apply exponents $(k,l)$ to $bd_0$
and $(k-1,l-1)$ to $abcd_0$. This gives \eqref{eq:u0}--\eqref{eq:u1}
and preserves simplicity and nonseparation.

For an additional handle let $p$ be the position of $a_i$ in the full
cyclic order. The closing chord $\{1,6\}$ of $abcd_0$ is replaced, for
$abcd_0a_i$, by $\{6,p\},\{1,p+2\}$, and for $abcd_0b_i^{-1}$ by
$\{6,p+1\},\{1,p+3\}$. These chords are again noninterleaving with each
other and with the three unchanged chords. Each edge is used once.
The resulting curves are simple and have nonzero homology classes.
The same boundary-fixing twists give the claimed words.
\end{proof}

We will also use the twist about $b$ whose based action is $a\mapsto ab$.
It fixes all other generators and satisfies $[ab,b]=[a,b]$. All twists
used below therefore lie in the orientation-preserving mapping class group
of the closed surface. No replacement of $ba^k$ by $a^kb$ is implicit.

\section{Cartan genericity and a residual linear system}\label{sec:linear}
Throughout this section $r\geq3$, and the eigenvalues $\mu_i$ are those
fixed in Lemma~\ref{lem:units}. All parameter spaces are over $\C$.

\begin{lemma}\label{lem:P}
There is a nonempty Zariski open set of $Q\in\SL_r$ on which both of the
following hold:
\begin{enumerate}[label=(P\arabic*),leftmargin=3em]
\item For every $i\ne j$, the map $\hc\to\C^2$,
$K\mapsto(K_{ij},K_{ji})$, has rank two.
\item The torus $\TC$ contains a regular matrix $N$ whose off-diagonal
entries are all nonzero and such that, for every three distinct indices,
\begin{equation}\label{eq:triangle-generic}
 N_{ij}N_{jk}N_{ki}\ne N_{ik}N_{kj}N_{ji}.
\end{equation}
\end{enumerate}
\end{lemma}
\begin{proof}
For (P1) at the pair $(1,2)$, take a matrix $H$ whose only nonzero
off-diagonal entries are
\[
 H_{12}=H_{21}=H_{23}=H_{31}=1;
\]
\emph{all other off-diagonal entries are zero}. Give it generic diagonal
entries. The projections of $H$ and $H^2$ to entries $(1,2),(2,1)$ are
\[
 (1,1),\qquad(H_{11}+H_{22},\ H_{11}+H_{22}+1),
\]
so their determinant is one. Generic diagonal entries make $H$ regular
semisimple (one may first take distinct widely separated real diagonal
entries). Subtract the scalar part of $H$; this does not change the
independence, since the change in $H^2$ is a scalar plus a multiple of $H$.
The two traceless polynomial expressions belong to the eigen-Cartan of
$H$. Choosing an eigenbasis and rescaling one column makes its matrix
$Q$ have determinant one. The rank-two condition is an open condition on
$Q$, since it is the nonvanishing of at least one $2\times2$ minor of a
fixed basis of $Q\h Q^{-1}$. Relabeling the witness treats every pair.
The finitely many nonempty opens intersect in the irreducible variety
$\SL_r$.

For (P2), the product of the off-diagonal coordinates, the finitely many
triangle differences in \eqref{eq:triangle-generic}, the determinant,
and the discriminant is a nonzero polynomial on the full matrix space:
each factor is a nonzero polynomial, and that polynomial ring is a domain.
Choose a matrix where it is nonzero and scale it to determinant one.
Its eigen-Cartan is a witness. More explicitly, substitute
$N=Q\diag(t_1,\ldots,t_r)Q^{-1}$, $\prod t_i=1$, in the product of the
required factors. This is a Laurent polynomial in $t_1,\ldots,t_{r-1}$
whose coefficients are regular in $Q$. Its not being identically zero is
a nonempty open condition on $Q$. It is equivalent to the existence of
a suitable $N$. Intersect with (P1).
\end{proof}

\begin{lemma}[Principal-minor differential rigidity]\label{lem:minors}
Suppose a matrix $M\in M_r(\C)$ has no zero off-diagonal entries and
satisfies \eqref{eq:triangle-generic} with $M$ in place of $N$. If a
variation $E\in M_r(\C)$ makes every principal minor stationary, then
$E=[h_0,M]$ for a diagonal matrix $h_0$.
\end{lemma}
\begin{proof}
The one-by-one minors give $E_{ii}=0$. Set $x_{ij}=E_{ij}/M_{ij}$ for
$i\ne j$. The two-by-two minors give $x_{ij}+x_{ji}=0$.
Differentiate the three-by-three determinant by the permutation formula.
Its diagonal term vanishes because $E_{ii}=0$; its transposition terms
vanish by the pair identities. The remaining two oriented cycles give
\[
 (M_{ij}M_{jk}M_{ki}-M_{ik}M_{kj}M_{ji})
                 (x_{ij}+x_{jk}+x_{ki})=0.
\]
The first factor is nonzero, so all oriented triangle sums vanish.
Set $h_1=0$ and $h_i=x_{i1}$ for $i\ne1$. Antisymmetry and the triangle
with reference index $1$ give $x_{ij}=h_i-h_j$ for every $i\ne j$,
including the cases involving index $1$. Thus
$E_{ij}=(h_i-h_j)M_{ij}$, including the diagonal entries.
\end{proof}

This is an infinitesimal form of classical principal-minor rigidity under
diagonal similarity. Loewy's global theorem \cite{Loewy} and the
pair-product and triangle-cycle identities of Mantelos
\cite[Sections 3--4 and Lemma 6.2]{Mantelos} are antecedents.
We give the elementary proof because its exact differential hypotheses
are what is needed below; no new contribution is claimed for this lemma
alone. In rank three the three-by-three minor is the full determinant;
in both applications its derivative vanishes automatically.

\begin{proposition}\label{prop:residual}
There is a nonempty Zariski open subset of $T\times\SL_r$ such that,
for $C=QC_\mu Q^{-1}$, the only $S\in\slr$ satisfying
\begin{equation}\label{eq:residual}
 (\Ad_A-\Ad_C)S=0,\qquad (I-\Ad_C)S\in\h+\hc
\end{equation}
is zero, and $\h\cap\hc=0$.
\end{proposition}
\begin{proof}
Choose a real orthogonal matrix $Q$ of determinant one with every entry
nonzero. Such matrices exist: each individual coordinate is nonzero on
an open dense subset of the real analytic manifold $\mathrm{SO}(r)$,
and a finite intersection is nonempty. Every member of $\hc$ is then a
complex symmetric matrix. Moreover $\h\cap\hc=0$: if
$\diag(k)Q=Q\diag(\lambda)$, nonzero $Q_{ij}$ implies
$k_i=\lambda_j$ for all $i,j$, so both diagonal matrices are scalar, and
tracelessness makes them zero.

Choose a real traceless diagonal $H$ with distinct entries, and put
$A(t)=\exp(tH)$ for real $t$. Suppose there were nonzero solutions of
\eqref{eq:residual} for a sequence $t\to0$, $t\ne0$. Normalize the
Hermitian norms of $S_t$ to one and pass to a convergent subsequence
$S_t\to K$. The first equation gives $K\in\hc$. It also gives
\[
 \frac{(I-\Ad_C)S_t}{t}
 =\frac{(I-\Ad_{A(t)})S_t}{t}\longrightarrow-[H,K].
\]
The left-hand side lies in the fixed closed subspace $\h+\hc$, hence
$[H,K]$ is complex symmetric. Since $H$ and $K$ are symmetric, their
commutator is skew-symmetric. Therefore $[H,K]=0$. Distinct entries of
$H$ make $K$ diagonal, whence $K\in\h\cap\hc=0$, contradicting its
norm one. Thus for all sufficiently small nonzero $t$ the system is
injective.

For global openness use the augmented linear map
\begin{equation}\label{eq:augmented}
\begin{split}
 \slr\oplus\h\oplus\h&\longrightarrow\slr\oplus\slr,\\
 (S,\alpha,\beta)&\longmapsto
 \bigl((\Ad_A-\Ad_C)S,\ (I-\Ad_C)S-\alpha-Q\beta Q^{-1}\bigr).
\end{split}
\end{equation}
It is injective exactly when \eqref{eq:residual} has no nonzero solution
and $\h\cap\hc=0$. Indeed a vector in its kernel gives the residual
system; after $S=0$, the remaining equality is exactly an intersection
of the two Cartans. Conversely each residual solution admits such
$\alpha,\beta$. All entries of \eqref{eq:augmented} are regular functions
on $T\times\SL_r$. Injectivity is a maximal-minor condition, and the
witness just constructed proves its nonemptiness.
\end{proof}

\begin{lemma}\label{lem:all-open}
There is a nonempty open set of triples $(A,B,Q)\in T\times T\times\SL_r$
satisfying (P1), (P2), Proposition~\ref{prop:residual}, and the following
additional conditions: $B$ has distinct diagonal entries; every entry of
$Q$ and $C=QC_\mu Q^{-1}$ is nonzero; and for each fixed cardinality
$1\leq j<r$, the products
\[
 a_I=\prod_{i\in I}a_i\qquad(|I|=j),\quad A=\diag(a_1,\ldots,a_r),
\]
are distinct for distinct subsets $I$.
\end{lemma}
\begin{proof}
Each condition is nonempty and open, and the parameter space is
irreducible. For the subset products, $\mathbf1_I-\mathbf1_J$ is not a
multiple of $(1,\ldots,1)$ when $I\ne J$ have equal size, so the forbidden
character equality is proper. Nonzero entries of $Q$ and regularity of
$B$ are elementary open conditions. Each off-diagonal entry of a matrix
in the regular conjugacy class of $C_\mu$ is nonzero somewhere: conjugate
a diagonal matrix by an elementary matrix and use two distinct eigenvalues.
Each diagonal entry is nonzero somewhere, since the eigenvalues themselves
are nonzero. Their individual nonempty opens therefore intersect.
The other conditions are covered by the preceding results.
\end{proof}

\section{The two-band cotangent certificate}\label{sec:tangents}
For a matrix $M\in\SL_r$ let $e_j(M)$ be the $j$th elementary symmetric
function of its eigenvalues, so
$\det(zI-M)=\sum_{j=0}^r(-1)^je_j(M)z^{r-j}$. We use $1\leq j<r$;
$e_r=1$ is constant. For a group word $w$, these are regular invariant
functions on the representation variety and descend to its character
variety.

The use of centralizers of twisted products and their powers has a close
precedent in Pickrell--Xia \cite[Proposition 2.1.24, Lemma 2.1.34]{PX}.
Here we need the following explicit calculation at a family whose first
two handle matrices commute; no generic compact-group spanning theorem
is substituted for the calculation.

\begin{proposition}\label{prop:core}
Let $(A,B,Q)$ satisfy Lemma~\ref{lem:all-open}, and set $C=QC_\mu Q^{-1}$.
At the genus-two representation $(A,B,C,I)$, the common tangent kernel of
the characteristic coefficients of $A,C$ and of all the simple curves
$u_0(k,l),u_1(k,l)$, $k,l\geq1$, consists exactly of infinitesimal
simultaneous conjugations. The resulting character is a smooth
irreducible point, and these differentials span its cotangent space.
\end{proposition}
\begin{proof}
Since $A,C$ are regular semisimple, stationarity of their characteristic
coefficients says that their variations are tangent to their conjugacy
classes. Subtract an infinitesimal simultaneous conjugation to write a
common-kernel tangent vector as
\begin{equation}\label{eq:initial-gauge}
 \delta A=0,\qquad \delta B=BU,\qquad
 \delta C=[Z,C],\qquad \delta D=V,
\end{equation}
with $U,V,Z\in\slr$. The choice of $Z$ is allowed to change by an element
of $\hc$.

\smallskip\noindent\emph{The first band.}
Cyclic rotation of a product of invertible matrices preserves its
characteristic polynomial. Use the cyclic representative $A^kDC^lB$
for $u_0$. Its derivative is $A^kE_l$, where
\begin{equation}\label{eq:first-variation}
 E_l=(V+Z)C^lB+C^l(\Ad_B U-Z)B.
\end{equation}
For any $M$ and any $j$,
\begin{equation}\label{eq:Vandermonde}
 e_j(A^kM)=\sum_{|I|=j}a_I^k\det M[I,I].
\end{equation}
For fixed $j<r$ the distinct nonzero numbers $a_I$ make the coefficient
matrix for $k=1,\ldots,\binom rj$ a nonsingular Vandermonde matrix times
a nonsingular diagonal matrix. Thus stationarity for all $k\geq1$
forces every $j\times j$ principal minor of $M=C^lB$ to be stationary
under $E_l$. The full determinant is stationary too: the variations come
from products in $\SL_r$.

Positive powers of $C$ are Zariski dense in $\TC$. To see that positive
powers suffice, their Zariski closure in the connected torus is the same
as the closure of the full cyclic group; equivalently a Laurent
polynomial vanishing on all positive powers also vanishes on all integer
powers by the recurrence satisfied by its finitely many exponential
terms. Therefore the resulting regular identities hold with $C^l$
replaced by every $N\in\TC$.
Put
\[
 X=V+Z,\qquad Y=\Ad_BU-Z,\qquad W=X+Y=V+\Ad_BU.
\]
All principal minors of $NB$ are stationary for the variation
$(XN+NY)B$. At $N=I$, the one-by-one identities give $W_{ii}=0$.
Differentiate the stationary two-by-two identity along $N=I+tK+O(t^2)$,
where $K\in\hc$. The one-by-one stationary identities hold for every
$N$, so the derivative of the diagonal-product part is zero. The two
off-diagonal terms give, after dividing by $-B_{ii}B_{jj}\ne0$,
\begin{equation}\label{eq:pair-W}
 K_{ij}W_{ji}+K_{ji}W_{ij}=0.
\end{equation}
Condition (P1) forces $W_{ij}=W_{ji}=0$ for each pair. Hence $W=0$ and
$V=-\Ad_BU$, and the variation reduces to $[X,N]B$.

Choose a regular $N$ as in (P2). The matrix $NB$ has nonzero
off-diagonal entries, and each triangle difference is that of $N$
multiplied by the nonzero scalar $B_{ii}B_{jj}B_{kk}$.
Lemma~\ref{lem:minors} gives
\[
 [X,N]B=[h_0,NB]=[h_0,N]B
\]
for a diagonal $h_0$. Thus $X-h_0$ centralizes regular $N$, whose
centralizer in $\SL_r$ is $\TC$. Subtracting scalar parts makes
$h_0$ traceless and $X-h_0\in\hc$. Since $X=Z-\Ad_BU$, subtracting
the simultaneous diagonal conjugation by $h_0$ yields precisely
\begin{equation}\label{eq:residual-gauge}
 \delta A=0,\qquad \delta B=SB,\qquad
 \delta C=[S,C],\qquad \delta D=-S,
 \qquad S=\Ad_BU.
\end{equation}
Here diagonal conjugation does not change $A$ or $B$, and the remaining
part of $Z-h_0-S$ commutes with $C$.

The derivative of the surface relation in \eqref{eq:residual-gauge} is
\[
 (\Ad_A S-S)+(S-\Ad_C S)=(\Ad_A-\Ad_C)S.
\]
Consequently any relation tangent in this gauge satisfies
\begin{equation}\label{eq:T0}
 T_0:=\Ad_A S=\Ad_C S.
\end{equation}

\smallskip\noindent\emph{The second band.}
At the base point $u_1(k,l)$ has matrix $M=A^kBC^l$.
Differentiate the actual word $AB A^{k-1}CD C^{l-1}$ in the gauge
\eqref{eq:residual-gauge}. The $B$ contribution is $T_0M$.
For the last three factors one has
\[
 \delta(CDC^{l-1})=[S,C^l]-CSC^{l-1}.
\]
Using $CS=T_0C$ gives
\begin{align}
 \delta M
 &=T_0M+A^kB(S-T_0)C^l-MS\notag\\
 &=[T_0,M]+A^kB[R,C^l],\qquad R=S-\Ad_CS.\label{eq:second-variation}
\end{align}
The conjugation term has zero characteristic derivatives. Apply
\eqref{eq:Vandermonde} and density of the powers of $C$ once more.
All principal minors of $BN$ are stationary for the variation $B[R,N]$,
for every $N\in\TC$. The full determinant is stationary because
$\tr(N^{-1}[R,N])=0$. At a matrix from (P2),
Lemma~\ref{lem:minors} yields
\[
 B[R,N]=[h_1,BN]=B[h_1,N],
\]
so $R-h_1\in\hc$ after removing scalar parts. Hence
$(I-\Ad_C)S\in\h+\hc$. Combined with \eqref{eq:T0}, this is exactly
\eqref{eq:residual}. Proposition~\ref{prop:residual} forces $S=0$.
Undoing the gauges proves the kernel assertion. The converse holds because
all the functions are conjugation invariant.

\smallskip\noindent\emph{Smoothness and the quotient.}
Since $B$ has distinct diagonal entries, a $B$-invariant subspace is a
sum of coordinate lines. No proper nonzero such sum is invariant under
$C$, because every entry of $C$ is nonzero. The representation is
irreducible. The image of the derivative of the relation contains
\[
 \im(\Ad_A-I)+\im(\Ad_C-I),
\]
by left variations of $B$ and variations of $D$. Under the nondegenerate
trace pairing its annihilator is the common Lie centralizer
$\h\cap\hc=0$. The relation map is therefore submersive. The relation
variety is smooth of dimension $3(r^2-1)$ at the point; simultaneous
conjugation has dimension $r^2-1$. Its stabilizer is the scalar center,
which acts trivially. The standard irreducible smooth quotient description
from Section~\ref{subsec:geometry} gives a smooth character point of
dimension $2(r^2-1)$, with tangent space the quotient by conjugations.
The common-kernel assertion is thus equivalent to cotangent spanning.
\end{proof}

\begin{proposition}\label{prop:all-genus}
For every $g\geq2$, at the representation
\begin{equation}\label{eq:family}
 \rho_A=(A,B,C,I,I,\ldots,I)
\end{equation}
with the same open conditions, the functions in
Proposition~\ref{prop:core}, together with the traces of
$u_1(k,l)a_i$ and $u_1(k,l)b_i^{-1}$ for $i\geq3$, span the full
cotangent space of $\X$. A finite list of these functions suffices.
\end{proposition}
\begin{proof}
The commutator of an extra pair of identity matrices has zero first
derivative, so its two tangent variables $\delta A_i,\delta B_i\in\slr$
are unrestricted by the linearized relation. Infinitesimal simultaneous
conjugation is zero on these identity entries. Apply the core argument
to a common-kernel tangent and gauge its first four variations to zero.
The trace differentials of the additional curves then are exactly
\begin{equation}\label{eq:extra-pairings}
 \tr(U_{kl}\delta A_i),\qquad-\tr(U_{kl}\delta B_i),
 \qquad U_{kl}=A^kBC^l.
\end{equation}
No subtraction of a core trace is necessary: its derivative is already
zero in this gauge.

The matrices $U_{kl}$ for $1\leq k,l\leq r$ span all of $M_r(\C)$.
Indeed $A,A^2,\ldots,A^r$ span the diagonal matrices by Vandermonde
interpolation. For fixed row $i$,
\[
 \operatorname{row}_i(C^l)=
   \sum_{j=1}^r\mu_j^lQ_{ij}\operatorname{row}_j(Q^{-1}).
\]
All $Q_{ij}$ are nonzero and the $\mu_j$ are distinct, so varying
$l=1,\ldots,r$ spans all row vectors. The additional nonzero factor
$B_{ii}$ does not change the span. The trace pairing against all of
$M_r$ detects every tangent matrix, so \eqref{eq:extra-pairings} kills
all extra variations. Smoothness follows by the same relation
submersion and irreducibility argument as before. The quotient dimension
is $(2g-2)(r^2-1)$.

A spanning set in a finite-dimensional cotangent space has a finite
spanning subset. Choose one such subset at one point in the nonempty
open parameter set, and let $\F$ be the finite set of underlying simple
curves. Retain the required coefficient index for each chosen function;
we may add any of the core functions used in constructing this list.
No effective bound on this finite list is asserted.
\end{proof}

\section{Arithmetic selection before complex nonresonance}\label{sec:selection}
The order of choices is essential. We first choose the finite list $\F$
from Proposition~\ref{prop:all-genus} at one suitable complex parameter
triple. We then choose integral $B,Q$ by finitely many open conditions.
Only after their choice do we impose nonresonance at a complex $A$.

\begin{lemma}\label{lem:finite-open}
For the fixed finite list $\F$, there is a nonempty open set
$\mathcal U\subset T\times\SL_r$ of parameters $(B,Q)$ with the following
properties:
\begin{enumerate}[label=(\roman*)]
\item On a nonempty open subset of the $A$-torus, $\rho_A$ is a smooth
irreducible character point and the selected differentials span its
cotangent space.
\item For each band curve in $\F$, its characteristic-coefficient map
from the $A$-torus to $T/W$ is dominant, where $W$ is the symmetric group.
\end{enumerate}
\end{lemma}
\begin{proof}
Put $q=r^2-1$. Use a fixed basis of the Lie algebra to trivialize the
tangent bundles of the matrix factors. At the initial suitable point, the
linear map consisting of the relation derivative together with the
selected invariant derivatives has kernel the $q$-dimensional conjugation
space. Its rank is therefore $2gq-q$. A corresponding minor is nonzero.
All entries are regular functions of $(A,B,Q)$: inverses are regular on
$\SL_r$, and all words are finite. If quotient charts were used instead,
one could clear their nonzero denominators with the same conclusion.
Include minors or products expressing relation submersion and the explicit
irreducibility conditions. This produces a nonzero regular function on
$T\times T\times\SL_r$ that guarantees (i). Expand it as a Laurent
polynomial in $A$. The locus of $(B,Q)$ where at least one coefficient
is nonzero is a nonempty open; there it is nonzero on a nonempty open
of the $A$-torus.

Along \eqref{eq:family}, each band in $\F$ has the same characteristic
polynomial as $A^kBC^l$ for some $k,l\geq1$. This includes $u_0$ by
cyclic rotation and commutation of $A$ with $B$; for the extra-handle
curves the extra matrix is the identity. At $Q=I$ the map in question is
\[
 T\longrightarrow T\longrightarrow T/W,\qquad
 A\longmapsto A^kBC_\mu^l\longmapsto(e_1,\ldots,e_{r-1}).
\]
The first arrow is finite surjective and the second is the finite Weyl
quotient. In characteristic zero its differential has full rank somewhere.
One such Jacobian minor is a Laurent polynomial in $A$, regular in $B,Q$.
The condition that it is not identically zero in $A$ is nonempty open
in $(B,Q)$ and implies dominance. This witness keeps exactly the fixed
eigenvalues $\mu_i$; no change of conjugacy class is being made.
There are only finitely many curves, so their opens and the open in (i)
intersect in the irreducible parameter space $T\times\SL_r$.
\end{proof}

Choose $(B,Q)\in\mathcal U$ with $B=B(m)$ and $Q\in\SL_r(\Z)$ by
Lemma~\ref{lem:arithmetic-dense}, and fix $C=QC_\mu Q^{-1}$. Then
\begin{equation}\label{eq:seed}
 \rho_0=(I,B,C,I,I,\ldots,I)
\end{equation}
is an actual representation into $\SL_r(\OO_L)$: every commutator is
identity. Let
\[
 Z=\overline{\Mod(\Sigma_g)\cdot[\rho_0]}^{\mathrm{Zar}}\subset\X.
\]
The twist $a\mapsto ab$ has powers sending $\rho_0$ to
$(B^n,B,C,I,\ldots,I)$. Since $B$ generates $T$ densely, continuity of
the quotient morphism gives
\begin{equation}\label{eq:torus-in-closure}
 [\rho_A]\in Z\qquad\text{for every }A\in T.
\end{equation}

\begin{lemma}\label{lem:choose-x}
There is an $A\in T$ such that $x=[\rho_A]$ is a smooth irreducible
character point where the selected differentials span the cotangent
space, and every curve in $\F$ has regular semisimple monodromy whose
powers are Zariski dense in its centralizer torus.
\end{lemma}
\begin{proof}
In a diagonal torus, failure of cyclic density is equivalent to
$\chi(t)=1$ for some nontrivial character
$\chi\in\Z^r/\Z(1,\ldots,1)$. For each such character its kernel is a
proper closed subset of $T$; the finite Weyl image of the union of its
conjugates is a proper closed subset of $T/W$. Properness follows from
dimension, since the quotient is finite. There are countably many such
characters. Also remove the discriminant locus in $T/W$ explicitly.
For each band map, dominance from Lemma~\ref{lem:finite-open} makes all
these inverse images proper closed subsets of the $A$-torus. For the
curve $a$ the spectral map of $A$ itself is dominant. For the curve $c$,
regular semisimplicity and nonresonance already follow from
Lemma~\ref{lem:units}, independently of $A$.

Intersect the nonempty open from Lemma~\ref{lem:finite-open}(i) with the
complements of this countable collection of proper closed sets, for all
curves in the finite list. Such an intersection is nonempty by the Baire
theorem: on any small closed ball contained in the open smooth torus,
proper complex algebraic subsets are closed with empty interior. The
selected $A$ has every required property, including avoidance of the
discriminant for each band.
\end{proof}

The $A$ of Lemma~\ref{lem:choose-x} is not asserted to be arithmetic.
This is harmless, because \eqref{eq:torus-in-closure} places it inside
the closure of the orbit of the single integral representation
\eqref{eq:seed}. Neither an arithmetic Baire argument nor a simultaneous
Hilbert-irreducibility assertion is used.

\section{Twist saturation and an open orbit closure}\label{sec:closure}
We complete the proof while allowing the orbit closure to be singular.
The torus saturation, Baire passage, and ambient submersion have
antecedents in \cite[Section on closed-surface dynamics]{Source027}.
Here all of them occur inside the closure of \emph{one} integral orbit.

Let $Y$ be the closure of the image of the $A$-torus in
\eqref{eq:torus-in-closure}. It is irreducible. Choose an irreducible
component $Z_0$ of $Z$ containing $Y$. For each $\gamma\in\F$, some
positive power $\tau_\gamma^{n_\gamma}$ of its Dehn twist preserves
$Z_0$, because each mapping class permutes the finitely many irreducible
components of $Z$.

\begin{lemma}\label{lem:component-nonres}
In a neighborhood of $x$ in $Z_0$, points whose monodromy on every
$\gamma\in\F$ generates its regular centralizer torus are dense in the
complex analytic topology.
\end{lemma}
\begin{proof}
Intersect $Z_0$ with the ambient smooth irreducible-character locus and
with the open where every $\gamma$-monodromy is regular semisimple.
This is a nonempty irreducible open $Z_0'$ containing $x$.
For one curve $\gamma$, pull back the eigenvalue-ordering cover
$T^{\mathrm{reg}}\to(T/W)^{\mathrm{reg}}$ to $Z_0'$. It is finite
\'etale. Each irreducible component of this cover dominates $Z_0'$:
flatness implies that its generic point lies over the generic point of
the integral base. Finiteness makes its image closed, so each component
in fact surjects onto $Z_0'$.

For any nontrivial character of $T$, its being one is a closed condition
on the cover. It cannot hold identically on any component, because every
component has a point over $x$, and every ordering of the eigenvalues at
$x$ is nonresonant. Its zero locus in each component is therefore proper;
under a finite morphism its image is closed and still proper, by dimension.
There are finitely many covering components and countably many characters.
Repeating this for the finite list of curves gives a countable collection
of proper closed subsets of $Z_0'$.

Their complement is dense analytically. For clarity, apply the Baire
theorem on small compact balls in the smooth dense locus of the
irreducible variety $Z_0'$, where each proper algebraic subset has empty
interior; then take closure to include singular points. This also shows
density in every relative analytic neighborhood of $x$.
\end{proof}

\begin{lemma}\label{lem:torus-orbit}
At each point $y$ in the dense set of Lemma~\ref{lem:component-nonres},
the centralizer-torus twist orbit about every $\gamma\in\F$ lies in $Z_0$.
\end{lemma}
\begin{proof}
Cut the surface along the simple curve $\gamma$. For a nonseparating
curve the surface group is an HNN extension of the cut-surface group.
Holding the cut representation fixed, one may multiply a transverse
gluing matrix by any element of the centralizer of $\rho_y(\gamma)$.
This is an algebraic torus action on the gluing parameter and a regular
map from that torus to the character variety. The Dehn twist is the
parameter $\rho_y(\gamma)$ or its inverse, according to orientation.
All curves in our list are nonseparating; the analogous separating
construction would conjugate one cut component.

The orbit under $\tau_\gamma^{n_\gamma}$ lies in $Z_0$. Powers of
$\rho_y(\gamma)^{n_\gamma}$ are still dense in its connected torus:
a character trivial on these powers would give a finite-order value on
$\rho_y(\gamma)$, whose suitable multiple would contradict its full
cyclic density. The inverse image of the closed set $Z_0$ under the torus
map is closed and contains these dense powers, so it is the whole torus.
\end{proof}

We next specify the complex Hamiltonian convention. On an ambient
neighborhood $U$ of $x$ contained in the smooth irreducible-character
locus $D/G$, let $\omega$ be the holomorphic symplectic form from the
complex trace pairing. For a holomorphic invariant function $h_\gamma$
let its holomorphic Hamiltonian field be defined by
\[
 \omega(X_{h_\gamma},\,\cdot\,)=\dd h_\gamma.
\]
The real two-form $\Omega=\Re\omega$ is nondegenerate: if
$\Re\omega(v,w)=0$ for all real tangent vectors $w$, also test $iw$ to
obtain the imaginary part. Under the identification of real tangent
vectors with their $(1,0)$ components, the real Hamiltonian of
$\Re h_\gamma$ has component $X_{h_\gamma}$, and the Hamiltonian of
$\Im h_\gamma$ has component $-iX_{h_\gamma}$.
Indeed these equalities follow directly from
$\Re(-i\dd h)=\Im\dd h$. Choosing $2\Re\omega$ and doubling the real
Hamiltonians gives the same fields.

This is also the variation relation of Goldman
\cite[Lemma 1.11]{Goldman}. His Theorem 4.3, proved in Theorems 4.5 and
4.7 (the HNN formula (4.6) in the nonseparating case), identifies these
real Hamiltonians with centralizer twists. The gradient of an invariant
function belongs to the centralizer Lie algebra. Therefore
$X_{h_\gamma}$ is tangent to the complex centralizer-torus orbit wherever
the monodromy is regular semisimple. This passage uses both real and
imaginary parts and takes place entirely in the ambient smooth locus;
it does not assume smoothness of $Z_0$.

\begin{lemma}\label{lem:ideal}
Every selected holomorphic Hamiltonian field has local complex-time flows
that preserve $Z_0\cap U$, including through its singular points.
\end{lemma}
\begin{proof}
Let $X$ be one such field and let $f$ be a holomorphic function on a small
open subset of $U$ vanishing on $Z_0$. At a point $y$ from
Lemma~\ref{lem:component-nonres}, Lemma~\ref{lem:torus-orbit} and the
Hamiltonian-twist identification give $(Xf)(y)=0$. By analytic continuity
and the density proved there, $Xf$ vanishes on all of $Z_0$ in that open
subset.

The radical ideal sheaf of the reduced analytic subspace $Z_0\cap U$ is
coherent. Near any point choose finite generators $f_1,\ldots,f_s$.
The preceding vanishing says that $Xf_j$ lies in that ideal sheaf; after
shrinking the neighborhood write
\[
 Xf_j=\sum_{k=1}^s a_{jk}f_k
\]
with holomorphic coefficients. Along an integral curve of $X$, the vector
$(f_1,\ldots,f_s)$ satisfies a homogeneous linear differential equation.
With zero initial values its unique solution is zero. Holomorphic ODE
existence and uniqueness, or the corresponding real ODE in each complex
time direction, shows that the local flow preserves the subspace even
when the starting point is singular on $Z_0$.
\end{proof}

\begin{proof}[Proof of Theorem~\ref{thm:main} for $r\geq3$]
At $x$, the chosen invariant differentials span the ambient cotangent
space by Lemma~\ref{lem:choose-x}. Nondegeneracy of $\omega$ makes their
Hamiltonian vectors span the ambient tangent space. Select a basis of
$d=(2g-2)(r^2-1)$ such vectors and compose their local holomorphic flows,
starting at $x$, with complex time parameters $(t_1,\ldots,t_d)$ near zero.
All intermediate points lie in $Z_0$ by Lemma~\ref{lem:ideal}. The
derivative of this map at the origin has the selected basis as its
columns, and is an isomorphism onto $T_x\X$. The holomorphic inverse
function theorem puts a nonempty ambient analytic open subset of $\X$
inside $Z_0$. Since $\X$ is irreducible by the exact source input in
Section~\ref{subsec:geometry}, such an open is Zariski dense. Thus
$Z_0=Z=\X$. The integral seed \eqref{eq:seed} proves the assertion.
The already credited rank-two case completes all $g,r$ in the statement.
\end{proof}

\appendix
\section{Proof scope and supporting artifacts}\label{sec:scope}
The proof of Theorem~\ref{thm:main} covers every $g\geq2$ and $r\geq2$
with the quantifiers stated there. Rank two is imported with explicit
credit; for $r\geq3$ Sections~\ref{sec:arithmetic}--\ref{sec:closure}
give the construction and proof. Classical surface-character geometry and
Goldman's Hamiltonian theorem are cited inputs. The irreducibility input
is the determinant-one intermediate assertion of the pinned manuscript,
identified in Section~\ref{subsec:geometry}. The proof neither classifies
all arithmetic orbits nor asserts analytic density of the selected orbit.

\subsection{Limited algebraic formal support}
A separate Lean artifact, \texttt{t027o-formal-v1}, supports a precisely
bounded part of Lemma~\ref{lem:minors}. Its theorem
\texttt{triangle\_coboundary}, over an arbitrary field and an index type
with a chosen reference index, proves that a zero-diagonal antisymmetric
array with vanishing oriented triangle sums is of the form
$x_{ij}=h_i-h_j$. Its theorem \texttt{principal\_minor\_rigidity} assumes
nonzero off-diagonal entries, zero diagonal variation, the explicit pair
stationarity polynomials, the six-term cycle derivative polynomials, and
unequal forward and reverse triangle products. It concludes
$E_{ij}=(h_i-h_j)M_{ij}$.

Those formal hypotheses are scalar polynomial identities. The pair identity
becomes the negative two-by-two minor derivative after the diagonal
identities; the cycle identity becomes the three-by-three minor derivative
only after the diagonal and pair identities. The artifact
does \emph{not} prove that they are the derivatives of the corresponding
determinant functions. It does not formalize generic Cartan choices,
Vandermonde interpolation, the surface curves, the residual degeneration,
arithmetic integrality, torus density, Goldman geometry, the analytic
ideal argument, or orbit density. It is not a formal proof of
Theorem~\ref{thm:main}. Its controls include a generic rational example
and a counterexample to deleting the triangle-genericity hypothesis.
This counterexample shows that genericity cannot simply be deleted; it does
not prove necessity of every triple condition in larger index sets. An
independent review adds a formal derivative bridge and further controls in
its own file, \texttt{T027OReview.lean}; those additions are distinct from
the author package and still do not cover the geometric orbit argument.
The exact coverage and trusted build boundary are documented in its
\texttt{t027o-formal-v1/README.md};
the manuscript packet supplies an external-path crosswalk for local use.
Formal review and release status are kept in the project records, not
inferred from the existence of this support file.

\subsection{Attribution and reproducibility limits}
The arithmetic unit construction, principal-minor cycle identities,
Vandermonde arguments, torus saturation, Baire theorem, coherent-ideal
argument, and inverse function theorem are supporting tools. They are not
separate novelty claims. The higher-rank contribution presented here is
the connected construction consisting of a single arithmetic seed,
the two-band tangent certificate at its complex torus family, and the
arithmetic-first, complex-second parameter choice. The literature
comparison in Section~\ref{sec:context} is bounded, with no claim of
complete priority clearance.

The source revision used is
\texttt{fd4aeeb2ee4fc729c18d98444fed42fd0529eeeb} of
\href{https://github.com/openai/math}{OpenAI's math repository}.
The author packet records exact source identities and reading limits.
Private copies of third-party bodies are excluded from the public source
bundle. The accompanying build script recompiles this article from a
copied source directory and records source and PDF hashes. These hashes
and finite diagnostic computations establish identity or finite checks;
they do not replace the mathematical proof.

\paragraph{Acknowledgment of assistance.}
The argument was developed with AI-assisted calculations, source
comparison, and independent project review. Author-side checks included
the unit-independence argument, the based-word conventions, the second-band
derivative, the residual symmetric/skew witness, and the spectral-cover
closure argument. Responsibility for the exposition and its stated
limitations remains with the named author. No institutional affiliation
is asserted.

\begin{thebibliography}{99}
\bibitem{BKMS}
I.~Biswas, T.~Koberda, M.~Mj, and R.~Santharoubane,
\emph{Representations of surface groups with finite mapping class group orbits},
New York J. Math. \textbf{24} (2018), 241--250.
\url{https://nyjm.albany.edu/j/2018/24-13p.pdf}.

\bibitem{CL}
S.~Coccia and D.~Litt,
\emph{Density of integral points in the Betti moduli of quasi-projective varieties},
arXiv:2507.00167v1 (2025).
\url{https://arxiv.org/abs/2507.00167}.

\bibitem{Goldman}
W.~M.~Goldman,
\emph{Invariant functions on Lie groups and Hamiltonian flows of surface group representations},
Invent. Math. \textbf{85} (1986), 263--302.
\url{https://www.math.umd.edu/~wmg/InvariantFunctions.pdf}.

\bibitem{GT}
A.~Salehi Golsefidy and N.~Tamam,
\emph{Closure of orbits of the pure mapping class group in the character variety},
Proc. Natl. Acad. Sci. USA \textbf{122} (2025), no.~15, e2416120122.
\url{https://doi.org/10.1073/pnas.2416120122}.

\bibitem{GTpre}
A.~Salehi Golsefidy and N.~Tamam,
\emph{Orbit closure of the pure mapping class group in the character variety},
detailed preprint, author-hosted version retrieved October 18, 2026.
\url{https://mathweb.ucsd.edu/~asalehig/G_T_OrbitClosure.pdf}.

\bibitem{Litt}
D.~Litt,
\emph{Motives, mapping class groups, and monodromy},
arXiv:2409.02234v1 (2024), especially Theorem~5.4.13.
\url{https://arxiv.org/abs/2409.02234}.

\bibitem{Loewy}
R.~Loewy,
\emph{Principal minors and diagonal similarity of matrices},
Linear Algebra Appl. \textbf{78} (1986), 23--64.
\url{https://doi.org/10.1016/0024-3795(86)90015-7}.

\bibitem{Mantelos}
H.~Sapranidis Mantelos,
\emph{Determinantally equivalent nonzero functions},
arXiv:2302.02471v4.
\url{https://arxiv.org/abs/2302.02471}.

\bibitem{MO}
C.~Frohman, question, with comments by M.~Kapovich and answers by
S.~Lawton and others,
\emph{Orbit of an irreducible representation of a surface group under that action
of the mapping class group}, MathOverflow question 332994 (2019 and later).
Discussion pointers, not an input theorem.
\url{https://mathoverflow.net/questions/332994}.

\begingroup\raggedright
\bibitem{Source027}
OpenAI,
\emph{Integral points on character varieties of curves},
September 25, 2026 manuscript, pinned repository revision
\texttt{fd4aeeb2ee4fc729c18d98444fed42fd0529eeeb}.
Directory \texttt{Integral-points-on-character-varieties-of-curves-September-25-2026}.
\url{https://github.com/openai/math}.\par
\endgroup

\bibitem{PX}
D.~Pickrell and E.~Z.~Xia,
\emph{Ergodicity of mapping class group actions on representation varieties,
I. Closed surfaces}, Comment. Math. Helv. \textbf{77} (2002), 339--362.
\url{https://ems.press/content/serial-article-files/42932}.

\bibitem{PreviteXia}
J.~P.~Previte and E.~Z.~Xia,
\emph{Topological dynamics on moduli spaces. II},
Trans. Amer. Math. Soc. \textbf{354} (2002), 2475--2494.
Theorem~1.4 is cited here only in the discussion of Litt's antecedent,
with the exceptional-locus caveat in \cite{GT}.
\end{thebibliography}

\end{document}
